
% Write about Core system, syntax overview, process structure, judgements, operational semantics etc. with focus on the multiplicative component.

% 3. πLL System
    % 3.1 Syntax (full system)
    % 3.2 Types (full system)
    % 3.3 Typing (full system)
    % 3.4 Operational Semantics (multiplicative fragment)

% Generally
    % Present syntax for processes and types
    % Show some Typing rules
    % For operational semantics, show TypingStep 
        % Define process and types semantics 
        % Make this very breif


\section{$\pi$MLL: Paper Level Definitions and Syntax}




% ------- WIP
\section{Operational Semantics}

\paragraph{Fragment decomposition.}
The system \pill can be naturally decomposed into two complementary fragments. The first is the multiplicative fragment, \pimll, which includes constructs such as tensor, par, and their associated units. This fragment captures the core structure of interaction, describing how processes communicate and synchronize through linear channels. The second is the non-multiplicative fragment, comprising additives, exponentials, and quantifiers. These constructs enrich the language with additional expressive power, enabling features such as internal and external choice, controlled duplication and disposal of resources, and polymorphic session behavior.

In this development, the syntactic and typing components are defined for the full \pill system, encompassing both fragments. However, the operational semantics and associated metatheory are established for the multiplicative fragment, which suffices to capture the essential communication behavior underlying the system. Extending the operational treatment to the non-multiplicative fragment introduces additional considerations, such as choice resolution and resource management, and is left for future work.
